\documentclass[12pt,a4paper]{article}

\setlength{\topmargin}{0.0in}
\setlength{\oddsidemargin}{0.33in}
\setlength{\textheight}{9.0in}
\setlength{\textwidth}{6.0in}
\renewcommand{\baselinestretch}{1.25}



\usepackage{hyperref}
% \usepackage[utf8]{inputenc}
% \usepackage[T1]{fontenc}
\usepackage{float}
\usepackage{enumitem}
\usepackage{amsmath}
\usepackage{amssymb}
\usepackage{tikz}
\usetikzlibrary{decorations.pathreplacing}
\usepackage{amsfonts}
\usepackage{graphicx}
\usepackage{subcaption}
\usepackage{caption}
\newcommand{\dv}[2]{\frac{\mathrm{d} #1}{\mathrm{d} #2}}
\usepackage{caption}
\captionsetup[figure]{font=small} 
% \captionsetup[table]{font=small}
\usepackage[capitalize,nameinlink]{cleveref}
\usepackage{mathtools}
\newcommand{\defeq}{\vcentcolon=}
\newcommand{\eqdef}{=\vcentcolon}

% Define environments
\newtheorem{theorem}{Theorem}
\newtheorem{definition}{Definition}
\newtheorem{lemma}{Lemma}
\newtheorem{corollary}{Corollary}



\title{Logs from June $2^{\text{nd}}$ to June $11^{\text{th}}$}
\author{}
\date{}
% \author{Maggie McCarthy}
% \date{May 2026}

\begin{document}

\maketitle

\section{Meeting June 2}

\subsection{Summary of last week's work:}




\begin{enumerate}
    \item Last week, I was focused on constructing the presentation.
    \item To construct the presentation, I performed a high-level study of Matt's algorithm and worked through the Maxwell example. I also had to consolidate some key spectral ideas, work through the shift operator examples, and clarify the project plan.
\end{enumerate}



\subsection{Meeting Notes}

\begin{itemize}
    \item In this meeting, I presented to Patrick, and he gave me comments slide-by-slide. I made note of these in a Google Doc.

    \item The next meeting is June 11, at 10:00 a.m. UK time.

    \item I have a meeting with Matt Colbrook sometime on June 11 (3-4pm).
    
\end{itemize}


\subsection{This Week's Plan}

\begin{enumerate}

    \item Fix up presentation and practice!
    \item Finish up some Boffi work.
    \item Transition to chapters 1 and 3 in Matt's text. 
    

\end{enumerate}


\section{Daily Logs}

\textbf{Tuesday} 
\begin{enumerate}
    \item Prepped for the meeting with Patrick, had the meeting.
    \item Spent the rest of the day making the desired changes to the presentation.  
\end{enumerate}

\textbf{Wednesday}
\begin{enumerate}
    \item More work on the presentation.
    \item Read through Matt's text to help my understanding before having to present it.
\end{enumerate}

\textbf{Thursday}
\begin{enumerate}
    \item Finalized up my slides
    \item Practiced all day! Fixed some issues as I went along.
    \item Made some presentation notes on what I could be asked. 
    \item Reviewed some of my background notes for the presentation.
\end{enumerate}

\textbf{Friday}
\begin{enumerate}
    \item Presentation morning - I think it went ok!
    \item Clarified the next steps with Patrick.
    \item Took an earlier day because my friend was visiting
\end{enumerate}


\textbf{Sunday}
\begin{enumerate}
\item Spent the morning reading about the de Rham complex and the Maxwell eigenproblem.
\item Finished up some typed of Laplace eigenvalue notes.
\item Started some Maxwell typed-up notes. I am hoping to finish these up either tomorrow or Tuesday morning. 
\end{enumerate}

\textbf{Monday}
\begin{enumerate}
\item Spent the day working through Part Two in Boffi's text. This took longer than expected, but it is a pretty dense section. 
\item Tomorrow I will finish up the last section of this part and start with Matt's text.
\end{enumerate}

\textbf{Tuesday}
\begin{enumerate}
\item Started the morning finishing up Part Two of Boffi. I also added some notes to my Maxwell example write-up. I want to work on these more in the future because I do not feel as though I fully understand why edge elements win!
\item Spent the rest of the day going through Forward, Preface, and the first half of Chapter 1 of Matt's text. The forward is beautiful, and may be helpful when I write the intro of the dissertation. Some examples in the first chapter are complicated - I tried to just take the key point of each example instead of diving into the properties of the operators presented. An interesting review of pseudospectra: Matt's definition differs slightly from Nick's (Matt uses closure of Nick's pseudospectrum). 
\end{enumerate}

\textbf{Wednesday}
\begin{enumerate}
    \item Started with a review of my FEM spectra notes, just to check for typos and clarity. 
    \item Worked through the rest of Chapter 1 of Matt's text. This took me almost the whole day which I did not expect.
    \item Finished by writing up some brief notes (bullet points) on Boffi's Part Two. 
\end{enumerate}

\end{document}
